\textbf{Culpability Without Ability: A Hard Determinist Proposal}

\textbf{Abstract}

This essay defends a coherent form of hard determinism against the
charge that, if agents are never able to do otherwise, our practices of
moral advice and blame lose their ground. I formulate this pressure as
the Culpability Dilemma. On the one hand, powerful incompatibilist
arguments deny any genuine ability to do otherwise, whether determinism
or indeterminism is true; on the other hand, ordinary moral life still
trades in judgments like `you should have...' and `it's your fault...',
judgments that seem to presuppose exactly that ability. To get out of
this dilemma, I challenge van Inwagen's Moral Advice Argument by
proposing that second-person `ought' claims are not primarily reports of
an antecedent normative fact that presupposes present categorical
ability but are practical interventions aimed at influencing what an
agent will do. On a deterministic picture, correct advice need only
satisfy a weaker condition: if the advice were given and taken up, it
would causally issue in the advised action. I then revisit Broad's
`ought implies can' by distinguishing comparative, deliberative, and
retrospective uses of `ought'. These revisions allow the hard
determinist to preserve the intelligibility of advice and much of the
grammar of fault and blame while rejecting the claim that culpability
requires the ability to do otherwise. The essential move is the
reinterpretation of Broad's retrospective ``ought'': showing that
past-tense moral verdicts reduce to counterfactual claims about correct
prior guidance and therefore do not presuppose categorical ability.

\textbf{1 Introduction}

We often blame people for certain states of affairs. We say a careless
driver is at fault for an accident, or that a dishonest politician is to
blame for a scandal. When we make such judgments, a trivial underlying
assumption seems to be that: there must be an occasion on which the
person we blame was able to do otherwise. We can more formally capture
this thought down in the following Culpability Thesis:

\begin{quote}
If a bad state of affairs S is truly the fault of a person x, then there
is some earlier time t and some course of action A such that: (1) At t,
x was choosing between two or more incompatible courses of action, one
of which is A; (2) at t, x ought to perform A instead of what x actually
did; (3) if x had performed A, S would not have occurred; and (4) at t,
x was able to perform A.
\end{quote}

The Culpability Thesis interacts in a crucial way with the following
three familiar philosophical theses:

\begin{quote}
Ability Thesis: On at least some occasions, when an agent is deciding
between two or more incompatible courses of action, the agent is able to
perform each of them.

Thesis of Determinism: The past and the laws of nature determine a
unique future.

Thesis of Indeterminism: The past and the laws of nature do not
determine a unique future.
\end{quote}

Both determinists and indeterminists can offer powerful arguments that
their respective views rule out the kind of ability posited by the
Ability Thesis. For the hard determinist (one who holds that the Thesis
of Determinism is true and that it is incompatible with the Ability
Thesis), Peter van Inwagen's Consequence Argument purports to show that
if the past and the laws of nature jointly determine every future event,
then no agent is ever able to do otherwise than what she actually does;
the ability to do otherwise would require either changing the past or
violating a law of nature, neither of which is within anyone's power.
For the hard indeterminist (one who holds that the Thesis of
Indeterminism is true and that it is incompatible with the Ability
Thesis), the Mind Argument powerfully shows that indeterminism fares no
better: if an agent's action is not determined by her prior mental
states, then whether she performs it is a matter of luck or chance, and
an action that occurs by chance is not one the agent has the ability to
perform or to refrain from performing in the relevant sense.\footnote{Peter
  van Inwagen, \emph{An Essay on Free Will}, rev. ed. (Oxford: Oxford
  University Press, 1986); and Peter van Inwagen, \emph{Thinking about
  Free Will} (Cambridge: Cambridge University Press, 2017), chap. 7,
  90--110, \url{https://doi.org/10.1017/9781316711101}.} Given that
determinism and indeterminism are mutually exclusive and jointly
exhaustive, if neither is compatible with the Ability Thesis, it follows
that the Ability Thesis is false. Yet the problem remains that in
ordinary life, we continue to speak and judge as if people are sometimes
at fault. When we think about Nazi officers who oversaw mass killings,
we say: ``These crimes were their fault'' or ``they ought not to have
done this,'' which, together with the Culpability Thesis, appears to
require exactly the sort of ability that both sides have just ruled out.

C. D. Broad found himself in a similar situation in his inaugural
lecture at Cambridge University, \emph{Determinism, Indeterminism, and
Libertarianism}, where he argued a libertarian view (one that affirms
the Ability Thesis) is hard to maintain given the incompatibilist
arguments. At this point, people like Broad, who deny the Ability Thesis
on the strength of these incompatibilist arguments face a dilemma with
two primary paths forward: (i) defend that ``nothing is ever anyone's
fault'' is not merely counterintuitive but actually true, and then
explain the nature of this powerful illusion that sustains ordinary
fault and blame attributions or (ii) argue that the Culpability Thesis
is false, i.e., moral responsibility does not require any ability to do
otherwise. This dilemma is what I will call the Culpability Dilemma.
Broad still believes the Culpability Thesis to be self-evident. His
reasoning is that whenever we make a retrospective judgment of
obligation (``you ought not to have done X''), we universally admit as
relevant the defense ``I could not help doing X'' or ``I could not
possibly have done Y.'' The fact that this is treated as a legitimate
rebuttal shows, for Broad, that there is some sense of ``could'' in
which ``ought'' and ``ought not'' entail ``could''; an action is, in his
words, ``obligable if and only if it is, in a certain sense,
substitutable.''\footnote{C. D. Broad, ``Determinism, Indeterminism, and
  Libertarianism,'' in \emph{Ethics and the History of Philosophy:
  Selected Essays} (London: Routledge, 2014), 197,
  \url{https://doi.org/10.4324/9781315823058}.} So (ii) is not really an
option for him.

Peter van Inwagen also builds his own case to defend the Culpability
Thesis and block path (ii) in \emph{Free Will and Moral Responsibility:
Once More Unto the Breach}, by connecting the concepts of fault, moral
advice, and the ability to act, in what I shall call the Moral Advice
Argument.\footnote{Peter van Inwagen, ``Free Will and Moral
  Responsibility: Once More unto the Breach,'' Belgrade Philosophical
  Annual 38, no. 1 (2025): 13, https://doi.org/10.5937/BPA2501007I.} His
argument is built on the core claim that: if ``you ought to do X'' is
the correct answer to the question ``What should I do?'', then the agent
is able to do X.

But does the Moral Advice Argument successfully establish the link
between culpability and ability or ``ought'' and ``can'', such that one
cannot coherently give advice or make fault and blame attributions if
they reject the Ability Thesis? This is the question I wish to explore
in what follows. Van Inwagen's core claim about advice, as I will show,
is a specific application of the Broadian principle: that ``ought'' and
``ought not'' can apply to an action only if the agent could have done
otherwise; so to establish culpability without ability requires showing
that the retrospective ``ought'' (the ``ought'' of blame and moral
verdict) does not commit us to the categorical sense that Broad takes to
be essential.

My strategy proceeds in stages. I will first (§2) offer a way around van
Inwagen's Moral Advice Argument by reconceiving advice as a causal
intervention, which weakens the link between correct advice and present
categorical ability. I will then (§3) turn to the deeper Broadian
principle itself and argue that once we distinguish the comparative,
deliberative, and retrospective uses of ``ought,'' the retrospective
``ought'' can be reinterpreted as a counterfactual claim about correct
prior guidance rather than as a marker of categorical ability. This
reinterpretation is the essential move to show that the hard determinist
can retain the full grammar of ``you ought to have done otherwise''
without accepting Broad's claim that such judgments presuppose the
ability to have done otherwise. Building on these two strands, I will
(§4) demonstrate how the hard determinist can make sense of fault and
blame ascriptions and dissolve the Culpability Dilemma via path (ii).

\textbf{2 The Moral Advice Argument}

Van Inwagen's Moral Advice Argument seeks to demonstrate that assigning
fault for a state of affairs logically implies that the agent in
question was, at some point, able to act otherwise. The core of the
argument can be formulated in three steps. Step one is a simple and
reasonable principle of fault, which states that:

\begin{quote}
``If a certain state of affairs is someone's fault, there is something
that person did not do and should have done (or did and should not have
done).''\footnote{van Inwagen, ``Free Will and Moral Responsibility,''
  15.}
\end{quote}

For example, if it's Henry's fault that the cat died, then Henry should
have fed the cat, or he should not have performed a cat-disabling
experiment (but he did). The object of blame is typically a state of
affairs (the cat's death), which is traced back to a failure to meet an
obligation (a thing Henry should have done). Step two is a trivially
true proposition:

\begin{quote}
``If an agent is deciding between A and B, and they ought to do A, then
the correct answer to the question `what should I do?' is `you ought to
do A.'\,''\footnote{Ibid.}
\end{quote}

This step simply blocks the objection that the ``ought'' of advice might
differ in content from the moral ``ought'': if there is a fact of the
matter about what one should do, then stating that fact constitutes
correct advice. When an advice columnist in ``The Ethicist'' gives
correct advice, the advice is correct precisely because it aligns with
what the person morally ought to do. Step three argues that:

\begin{quote}
``If `you ought to do x' is the correct answer to an agent's
deliberative question `what should I do?', then the agent is able to do
x.''\footnote{Ibid., 16}
\end{quote}

This is the critical step that establishes the `ought implies can'
principle. This step is, in effect, a particular application of a
principle that Broad articulated in his treatment of obligability: that
``ought'' and ``ought not'' can only apply to an action that is, in the
relevant sense, substitutable -\/- one that the agent could have
performed or could have left undone.\footnote{Broad, ``Determinism,
  Indeterminism, and Libertarianism,'' 195-197.} Van Inwagen's Step 3
applies this principle to the specific context of advice. The
disagreement in what follows will therefore concern not only van
Inwagen's particular formulation, but the deeper Broadian commitment
that stands behind it, which I will elaborate on in §3.

To illustrate why advising someone to perform an action that is beyond
their ability is senseless, van Inwagen used the following example:

\begin{quote}
Suppose Planner is giving his friend advice on where to take Alison for
dinner. Planner advises his friend to take Alison to Auberge du Vieux
Puits in Paris, despite knowing full well that his friend cannot afford
the trip and could not get there in time.
\end{quote}

To see how each steps play together, we can apply them to the Henry
feeding the cat example again.

1. Suppose the cat's death is Henry's fault (starting premise).

2. By step one, Henry should have fed the cat.

\begin{quote}
3. By step two, if Henry deliberates whether to feed the cat or go to
Camp David, the correct answer to ``What should l do?'' is: You ought to
feed the cat.

4. By step three, for that advice to be correct, Henry must have been
able to feed the cat.
\end{quote}

5. Therefore, if the cat's death was Henry's fault, Henry was able to
feed the cat.

\textbf{2.1 Moral Advice as Attempt to Influence}

The problem of this argument lies in step three which seems to
presuppose a hidden principle:

\begin{quote}
(CA.α) In order for X to give Y the correct advice, X must believe that
Y is able to do the thing advised.
\end{quote}

Planner's advice, ``you ought to take Alison to Auberge du Vieux Puits
for dinner'' is a wrong answer to his friend's question ``where should I
take her to dinner?'', because Planner does not believe that his friend
is able to do the thing advised, violating CA.α.

If we build CA.α into step three, we get:

\begin{quote}
Strong Step Three (SS.3): If ``You ought to do x'' is the correct answer
to an agent's deliberative question ``What should I do?'', then the
agent is (at the time of advice) able to do x.
\end{quote}

Following SS.3, van Inwagen's point is that anyone who rejects the
Ability Thesis simply cannot give out any moral advice that is
consistent with their belief in that rejection; to the question ``what
should I do?'', they cannot correctly answer ``you ought to do x.'' To
illustrate his reasoning, suppose that person X is a hard determinist.
If Y asks X: ``What should I do, A or B?'', from X's epistemically
limited perspective, X knows that Y will only do what Y does; since A
and B mutually exclusive, X must believe that Y is not able to do at
least one of them, but X does not know which. Therefore, on that ground,
X would have to say to Y: ``I cannot give you any advice, because I
don't know what you're able to do.''

However, this is not necessarily the case. In this scenario, what X
could say is that: ``once I gave Y the advice, then, Y would be able to
do the thing because he would follow my advice, so it's worth a try.''
To elucidate, it helps to distinguish two senses in which giving advice
might be ``worth a try.'' At the metaphysical level, under hard
determinism, whether X's advice will causally influence Y is already
settled: the prior state of the world and the laws of nature determine
whether the advice will be taken up, just as they determine everything
else. In this sense, it is not ``worth a try'' -\/- the outcome is
fixed. But at the epistemic level, X does not know what Y is determined
to do. X does not know whether the advice will be part of the causal
chain that leads Y to act, or whether it will fall on deaf ears. From
X's epistemically limited perspective, giving advice is worth a try in
the sense that, for all X knows, the advice might be the very causal
factor that determines Y's action.

This reframes what X can say when Y asks, ``What should I do, A or B?''
X need not claim to know that Y is presently able to do either. Rather,
X can reason as follows: ``I do not know, at this moment, whether Y will
be able to do what I advise. But I do know that, if I give the advice
and it is taken up, it will causally contribute to Y's doing the advised
action.'' Since X cannot know in advance whether the advice will be
efficacious, the rational course is to advise as if it might be - to
treat the giving of advice as a causal intervention whose success is
epistemically open, even if metaphysically settled.

The ``able to do A'' condition need not hold at the time of asking but
can be satisfied at the time of action after the advice. In this sense,
we can now replace CA.α with a more nuanced principle:

\begin{quote}
(CA.β) In order for X to give Y the correct advice, X has to believe
that Y is able to do the thing advised after receiving the advice.
\end{quote}

Going back to step three, notice that the strong version secretly builds
in the requirement of the agent's ability at the moment of advice. It
states you may give it as correct advice only if the agent,
independently of your giving it, already has the power to comply. A hard
determinist can reject this as the wrong picture of advice. In a
deterministic picture, when Y comes to X and asks, ``What should I do, A
or B?'', X is not standing outside the causal order trying to read off,
like a fortune teller, whose actions are already in Y's ability set. X
is trying to influence Y's decision by adding a cause or a determining
factor, namely, the advice.

Under this conception of advice, a hard determinist can then build CA.β
into step three and have a weaker version:

\begin{quote}
Weak Step Three (WS.3): If ``You ought to do x'' is the correct answer
to an agent's deliberative question ``What should I do?'', then the
agent is (after the advice) able to do x when the time for action comes.
\end{quote}

We can see how the revised WS.3 can change the Moral Advice Argument by
applying it to the Henry's cat case. Steps 1-3 remains the same, but the
revised steps 4 and 5 now reads:

\begin{quote}
4. If ``You ought to feed the cat'' is the correct answer, then Henry
must have been able to feed the cat (at the time of action) as a direct
result of the advice itself.

5. Therefore, if the cat's death was Henry's fault, then there is a
correct piece of advice such that, had it been given, Henry would have
been able to feed the cat at the time of action.
\end{quote}

The cat's death being Henry's fault only proves that the correct advice
``feed the cat'' would have been a causally efficacious factor in
bringing about that action. It does not prove that Henry had any ability
to choose that path independently of being caused to do so by the
advice.

Those are different modal claims. SS.3 licenses an inference from
``correct advice'' to alternative possibility: holding fixed the actual
past and the laws, Henry had both options genuinely open. WS.3 only
licenses an inference from ``correct advice'' to counterfactual,
advice-enabled possibility: in a nearby possible world where the past
includes the advice, Henry feeds the cat. A natural objection is: what
if the advice fails? What if Henry receives the advice and, being
determined by his prior psychological states and circumstances, does not
take it up? This is a real possibility, and the hard determinist should
concede it since WS.3 does not guarantee that advice will always be
efficacious. What it guarantees is weaker: that the correctness of
advice does not presuppose any ability on the agent's part at the time
of asking. Whether the advice succeeds is a further empirical question,
one whose answer is metaphysically determined but epistemically open to
the adviser. The hard determinist gives advice under the same epistemic
conditions as anyone else - not knowing for certain whether it will
work, but having reason to think it might, and recognizing that the
giving of advice is itself a potential causal factor in the outcome.

\textbf{2.2 Moral Advice and Non-Cognitivism}

The account of advice I just sketched relies on a simple thought: when
someone asks, ``What should I do, A or B?'', they are not merely asking
for a description of how the moral universe already stands; they are
asking for guidance that is meant to influence what they will do.

This way of thinking sits comfortably with many forms of non-cognitive
or moral subjective views of morality. Suppose a person thinks that
there are moral truths about things, then moral advice is just supposed
to be a report about what the moral universe says. Asking, ``what should
I do, A or B?'' is like asking, ``which number is bigger, 7 to the 3 or
3 to the 7?'' There's an answer to this question, and answering it is to
report the truth, not to guide someone's arithmetic. That's the point of
cognitivist morality, that my account of moral advice is inconsistent
with.

This non-cognitive aspect is a useful note to register explicitly how my
conception of advice differs from van Inwagen's. In his view,
second-person `ought'-statements are primarily informative or
reportarial. When an agent asks, ``What should I do, A or B?'', van
Inwagen understands this as a request for a judgment about an objective
normative fact; the adviser's role is to provide an answer that
accurately reports how things stand: what the right act is. What the
advisee then does with that information is, strictly speaking, up to
them and lies outside the content of the advice itself. In this picture,
the adviser does not ``have a dog in the fight'' between the agent's
competing inclinations; the adviser is not trying to strengthen one
desire or create a new one, but to tell the truth about which option is
right. By contrast, the hard-determinist reinterpretation treats
second-person ``ought''-claims as practical interventions in a causal
process. When X says to Y, ``ABC is the best restaurant,'' the primary
role of this utterance is not to neutrally report a normative fact for Y
to do with as Y pleases, but to strengthen or create an inclination in Y
and thereby help determine what Y will in fact do. The content of the
utterance is inseparable from its role in trying to steer behavior down
a certain path. This is not to say that the advice lacks evaluable
content: X can still be wrong about which restaurant is best, and if X
is wrong, then X has given bad advice. The point is that the function of
the utterance in the deliberative context is interventionist rather than
reportorial. Whether the advice is correct (by whatever evaluative
standards one endorses) and whether it is causally efficacious are two
separate questions. The hard determinist can accept that advice is
assessable for correctness while maintaining that its role in the
practice of advising is to intervene in a causal process whose outcome
is metaphysically determined but epistemically open to both adviser and
advisee.

The Moral Advice Argument implicitly presupposes the more
``objective-report'' conception. My proposal is that a hard determinist
can instead adopt the ``interventionist'' conception, and once that
shift is made, the move from ``correct advice'' to present categorical
ability in SS.3 no longer follows.

\textbf{2.3 Moral Advice and Hard Indeterminism}

Another important note is that this conception of moral advice does not
work for the hard indeterminists and only works under a specific
conception of hard determinism.

To see why, it helps to state more precisely what CA.β requires. On the
closest-world semantics for counterfactuals, CA.β is satisfied when the
following nested counterfactual is true:

\begin{quote}
Had the agent been advised to do A, then, had the agent tried to do A,
the agent would have done A.
\end{quote}

The outer counterfactual shifts us to the nearest possible world in
which the agent receives the advice; the inner counterfactual, evaluated
at that world, asks whether trying to comply would have resulted in
compliance. Crucially, this is a claim about what counterfactuals are in
fact true - a matter of the modal structure of the world - and not a
claim about what X knows. X gives advice under epistemic uncertainty
about whether these counterfactuals hold, just as a doctor prescribes a
treatment without certainty that it will work. What matters is whether
the counterfactuals are true, not whether X can verify them.

With this in hand, we can see why the same proposal fails for the hard
indeterminist. Suppose X is a hard indeterminist. If Y asks X, ``What
should I do, A or B?'', X does not know, prior to giving advice, what Y
is able to do. But crucially, even after giving the advice, X cannot be
sure that Y will be able to do the action advised - and this is not
merely an epistemic limitation. In X's picture of the world, the
relevant counterfactual may lack a determinate truth-value: even in the
closest possible world in which Y receives and tries to follow the
advice, indeterministic factors may prevent compliance. The problem is
the metaphysical level, not epistemic. Under determinism, if the advice
enters the causal history, the closest world in which Y receives and
tries to follow it is one in which Y succeeds (given causal stability).
Under indeterminism, no such assurance is available even at the level of
counterfactual truth, because chance may intervene between trying and
succeeding.

This contrast also reveals a tacit presupposition in the deterministic
case. My account of advice assumes a deterministic world with relatively
low modal fragility at the level of ordinary deliberation: small
variations in how advice is delivered do not produce wildly divergent
outcomes within the timeframes relevant to deliberation. This is a claim
about the modal structure of the world, not about what the adviser
knows. It does not require infallible modal knowledge on X's part; it
requires only that, as a matter of fact, the closest worlds in which the
same advice is given with minor variations (delivered in 19 seconds
rather than 20 or phrased slightly differently) are worlds in which the
same outcome obtains. If the world were so chaotic that this were not
the case, then the relevant counterfactuals would be false, and CA.β
would not be satisfied - but this would be a failure of the world to
cooperate, not a failure of X's knowledge.

My argument therefore supposes that, within the timeframes and levels of
description relevant to ordinary deliberation, advice functions as a
reasonably stable causal influence. This supposition goes beyond
standard hard determinism, but it is precisely the sort of stability
that our everyday advisory practices already take for granted.

\textbf{3 Does `Ought' Imply `Can'?}

So far, I have focused on van Inwagen's Moral Advice Argument and how a
hard determinist can reinterpret the link between advice and ability.
The deeper question remains: how can the hard determinist coherently use
the language of `ought' at all, or in other words, does all `ought'
imply `can'? Having loosened van Inwagen's particular formulation, in
this section, I turn to Broad's treatment of this deeper question in his
inaugural lecture, \emph{Determinism, Indeterminism, and
Libertarianism}. I will argue that Broad's conception of `obligability'
is tied to a particular retrospective use of `ought', and he does not
clearly distinguish this from other, forward-looking uses of `ought' as
advice. On that basis, I reconstruct his two senses of `ought' in a way
that a hard determinist can retain.

\textbf{3.1 Broad on `Ought Implies Can'}

Broad starts from a familiar kind of retrospective moral judgment about
a past action. We say things like: ``You ought not to have done A; you
ought instead to have done B.'' He then notes that if the person wants
to resist this judgment, they have two options. Either they say: (a) ``I
could have done A instead of B, but you are wrong about what I ought to
have done; in fact, A was what I ought to have done,'' or they say: (b)
``I could not help doing A,'' or ``I could not possibly have done B.''
In (a), the agent grants that `ought' or `ought not' applies, but
disputes which action is the one they ought to have done. In (b), the
agent denies that `ought' or `ought not' applies at all, because he or
she was not able to do otherwise. Broad takes (b) to rest on something
like a shared background assumption: if an agent literally was not able
to do otherwise, then neither `ought' nor `ought not' applies to that
action. This is what leads him to the thesis that `obligability entails
substitutability'.

In Broad's terminology, an action is ``obligable'' if it is the kind of
thing to which `ought' or `ought not' can be applied. An action is
`substitutable'' if it was done but was able to have been left undone,
or it was left undone but was able to have been done.\footnote{Broad,
  ``Determinism, Indeterminism, and Libertarianism.'' In the actual
  paper, Broad used the classical phrase ``could have'' as opposed to
  ``was able to''. However, given that ``could have'' is ambiguous as
  mentioned in 1.1, I assume what he meant by ``could have'' is really
  just ``was able to'', as opposed to ``might have''.} He then tries to
make this more precise by clarifying that the kind of substitutability
entailed by obligability is a `categorical substitutability', one that
does not share the logical form: ``I would have been able to do
otherwise if I had willed otherwise'' or ``I would have willed otherwise
if I had willed differently in the past.'' Broad's view can hence be
summarized in the following thesis:

\begin{quote}
``\,`Ought' or `ought not' can only apply to an action if the action was
done but was able to have been left undone, or it was left undone but
was able to have been done.''\footnote{Broad, ``Determinism,
  Indeterminism, and Libertarianism,'' 197.}
\end{quote}

Broad also holds that there are two general senses of `ought'. The first
he attributes to Spinoza and calls the `comparative ought'; This `ought'
can be applied to both humans and non-humans alike which simply states
that one state of affairs ranks higher than its competitors. For
example, suppose a person says, ``This soup ought to have less salt.'';
what the person meant is that its flavor is worse than that of a
properly seasoned soup. Similarly, when applied to human conduct, ``You
ought not to cheat at poker'' may mean that cheating falls below the
standard of an average decent person. Such claims are entirely
independent of whether the agent is able to meet that standard; this
sense, when applied to humans, is perfectly consistent with an absence
of categorical substitutability or the denial of the Ability Thesis.

The second sense of `ought' is what Broad called, the `categorical
ought'. When a person feels guilty and thinks, ``I ought not to have
done that,'' the person is not just saying, ``my character is worse than
average''; the person is implying that he or she was able to have and
should have chosen differently, independent of his or her pre-existing
character. This is the sense of `ought' that implies categorical ability
and hence making the Culpability Thesis more believable.

\textbf{3.2 The Grammars of `Ought'}

One important note for Broad's arguments is that he never distinguishes
between the retrospective `ought' and the moral advice `ought'. In fact,
Broad's examples and formulations in the lecture almost all concern
past-tense judgments:

``You ought not to have done X.'',

\begin{quote}
``You ought to have done Y instead.'',

``He ought not to have let himself become a slave to morphine,'' etc.
\end{quote}

For Broad, the grammatical form here is that we look back at what
someone did (or failed to do) and pronounce a verdict. Broad calls these
retrospective judgments of obligation about past actions, and his whole
discussion of `obligability' and `substitutability' is framed in terms
of such judgments.

What he does not do is to distinguish this retrospective `ought' from
other very common uses of the word we discussed in §2, namely: the
future-directed `ought' of moral advice (``You ought to tell the truth
tomorrow''). Instead, Broad slides from his careful discussion of
past-tense `ought' into a more general-sounding principle of `ought
implies can'. It is then natural for readers to take this as a
constraint on all uses of `ought'.

But this slide can be resisted. We can distinguish between different
roles that `ought' plays in our language and re-interpret this
retrospective `ought' in a way that aligns with the future-directed
`ought'.

Drawing on the earlier discussion of moral advice, we can now
distinguish at least three useful senses or roles of `ought':

(i) First, we still have the `comparative ought'. This Spinozan sense of
`ought', as sketched out by Broad, is certainly compatible with hard
determinism.

(ii) Second, we have the `deliberative ought'. This is the sense that
appears when X asks, ``What should I do, A or B?'' and Y answers, ``You
ought to do A.'' Given the picture developed in §2, this `ought' is tied
to the practice of advice as an attempt to influence or guide action and
hence can also be used consistently under hard determinism.

(iii) Third, we have the `retrospective ought'. This is the past-tense
`ought' that Broad focuses on when one looks back and makes a normative
verdict about a past action. Based on Broad's view, this is a
`categorical ought' that is rendered false if the agent was not able to
do otherwise.

The problem is how can hard determinists interpret the `retrospective
ought' in a way that does not commit them to the categorical sense. We
can use the insight from WS.3 in §2; a hard determinist can analyze
retrospective judgments of the form ``X ought (not) to have done ϕ'' by
reducing them to counterfactual claims about correct prior advice. More
precisely:

(R) ``X ought (not) to have done ϕ at t'' ⇔ ``If, at some suitable time
t*, prior to t, X had asked `What should I do?', then the correct answer
would have been: `You ought (not) to do ϕ.'\,''

For a hard determinist, to make a retrospective judgment and say ``you
ought to have done or ought not to have done something'' is like saying,
``the advice I would have given you if you'd asked beforehand is x.'' On
this analysis, a retrospective `ought' is a claim about what the correct
deliberative `ought' would have been, had X sought advice or guidance
beforehand. The truth-conditions of the retrospective ``You ought to
have done ϕ'' are therefore identical to the truth-conditions of the
corresponding advisory claim ``You ought to do ϕ'' evaluated at an
earlier time t*. In this way, the retrospective `ought' inherits the
interpretation of the deliberative `ought' the hard determinist adopts
in CA.β, without introducing any further commitment to categorical
ability.

Therefore, when X says to Y, ``you ought to have done A or ought not to
have done B,'' X might be doing at least two things:

\begin{enumerate}
\def\labelenumi{(\arabic{enumi})}
\item
  Evaluating Y's past action relative to a standard (``Y's action falls
  below the standard for a decent person''), which is a comparative
  `ought',
\end{enumerate}

\begin{quote}
or
\end{quote}

\begin{enumerate}
\def\labelenumi{(\arabic{enumi})}
\setcounter{enumi}{1}
\item
  Indicating what the correct advice would have been and,
  simultaneously, issuing a forward-looking recommendation (``In similar
  situations, you ought to do A and not do B''), which is effectively a
  deliberative `ought'.
\end{enumerate}

Under this interpretation, there are still essentially two senses of
`ought', both of which are compatible with hard determinism, undermining
the `ought implies can' principle as argued by Broad.

\textbf{4 Out of the Culpability Dilemma}

So far, I have argued that a hard determinist can make sense of
second-person moral advice by treating ``You ought to do A'' as a causal
intervention rather than as a neutral report about an independent
normative fact. A natural next question is whether the same strategy can
be extended to fault and blame.

The hard determinist can answer this by extending the same causal
interventionist account from advice to blame. To say ``It is your fault
that S happened'' has at least two components. First, there is a
descriptive, cognitive component: it attributes to the agent a certain
causal and normative relation to S; it notes that an action or omission
by the agent played a role in the deterministic history leading to S and
evaluate that action as falling below some accepted moral standard. In
this sense, the hard determinist can say ``You are at fault,'' to mean
``Your behavior was one of the causes of S and manifested a defect
relative to the standard we endorse.'' No appeal to categorical
abilities is required for that much.

Secondly, there is a forward-looking, non-cognitive component. In
ordinary life, we use ``You are to blame'' to produce feelings of guilt
and remorse to shape the agent's future disposition or to reinforce a
shared norm in the community. In the interventionist picture,
ascriptions of fault and blame are themselves attempts to influence the
future causal chain in much the way that advice is.

This allows the hard determinist to reasonably reject the ability
requirement found in clause (3) of the Culpability Thesis and instead
substitute a weaker condition, parallel to the Weak Step Three from
§2.1. Accordingly, we can construct a hard deterministic reformulation
of the Culpability Thesis as follows:

\begin{quote}
A bad state of affairs S is truly the fault of a person x if and only if
(1*) x's action or omission was a part of the causal history of S, (2*)
that action or omission falls below a standard that we are entitled to
enforce, and (3*) there is some correct pattern of guidance or advice
such that, had x been appropriately exposed to and responsive to it
earlier, x would have behaved differently and S would not have occurred.
\end{quote}

On this view, fault-ascriptions, like advice, are part of the machinery
by which a deterministic system regulates itself.

We can now see how this returns us to the Culpability Dilemma. Given
that the practice of any fault and blame ascriptions require only causal
connection to bad outcomes, deviation from standards we are prepared to
enforce, and a forward-looking role for blame and advice as instruments
for influencing behavior, we can consistently deny `ought implies can'
under hard determinism, and path (ii) remains open to get out of the
Culpability Dilemma.

\textbf{Bibliography}

Broad, C. D. ``Determinism, Indeterminism, and Libertarianism.''
\emph{Ethics and the History of Philosophy: Selected Essays}. Routledge,
2014. \url{https://doi.org/10.4324/9781315823058}.

van \href{https://www.zotero.org/google-docs/?QMZQDU}{Inwagen, Peter.
\emph{An Essay on Free Will}.} Rev. ed. Oxford: Oxford University Press,
1986.

van Inwagen, Peter. \emph{Thinking about Free Will}. Cambridge:
Cambridge University Press, 2017.
\href{https://doi.org/10.1017/9781316711101}{{https://doi.org/10.1017/9781316711101}}.

van Inwagen, Peter. ``Free Will and Moral Responsibility: Once More unto
the Breach.'' \emph{Belgrade Philosophical Annual} 38, no. 1 (2025):
7--19.
\href{https://doi.org/10.5937/BPA2501007I}{{https://doi.org/10.5937/BPA2501007I}}.

van Inwagen, Peter. ``Does `Ought' Imply `Can'?'' \emph{Free Will:
Historical and Analytic Perspectives}, edited by Marco Hausmann and Jörg
Noller, 313--33. Cham: Springer, 2021.
